We introduce our approach as feature-composed I-ProtoMF, in short
fI-ProtoMF. It differs from I-ProtoMF~\cite{melchiorre2022protomf} in
exactly one place. The host draws the embedding $\mathbf{t}$ of every
item from a lookup table with one free learned row per item. fI-ProtoMF
replaces this table by a composition in the manner of
LightFM~\cite{kula2015lightfm},
\begin{equation}
\mathbf{t} \,=\, \sum_{f \in \mathcal{F}_t} \mathbf{e}_f
\;\,\bigl(+\; \mathbf{e}_{\mathrm{ID}(t)}\bigr),
\label{eq:fi-composition}
\end{equation}
where $\mathcal{F}_t$ holds the feature values recorded for the item,
each $\mathbf{e}_f \in \mathbb{R}^d$ is a learned vector, and
$\mathbf{e}_{\mathrm{ID}(t)}$ is an optional per-item row discussed
below. Everything downstream remains the host. The $L^t$ item prototypes
$\mathbf{p}^t_l$ stay free learned vectors, the item representation
$\mathbf{t}^* \in \mathbb{R}^{L^t}$ is still the vector of similarities
$\mathrm{sim}(\mathbf{t}, \mathbf{p}^t_l)$, the user is still a free
vector $\mathbf{u} \in \mathbb{R}^{L^t}$, and the score
$\text{I-score}(u,t) = \mathbf{u}^\top \mathbf{t}^*$ is trained with the
host's loss and regularizers (Section~\ref{sec:bg-protomf}). The model
name keeps the f of feature rather than attribute because the
composition sums whatever encoded rows it is given, the ID row being the
first example, and choosing attribute-derived rows is what makes the
resulting model interpretable.

The composition turns attribute values into learned words. Every
distinct value that appears in the catalogue, the colour group black or
the garment group jersey, owns exactly one row $\mathbf{e}_f$, shared by
all items that carry it. An item is then spelled from these words. A
plain black t-shirt sums the rows of its recorded values for department,
product type, section, colour group, and pattern, and a second shirt
that differs only in colour shares four of its five summands. The host's
item table, one independent row per item, becomes a small vocabulary of
attribute rows from which all item embeddings are assembled. The rows
are trained by the collaborative signal like any other parameter, but
each row receives the gradients of every interaction with every item
that carries its value, so what a row comes to encode is the
collaborative meaning of one attribute value across the catalogue.

This is the grounding announced in the previous section, stated
precisely. What the composition constrains directly is the item factor.
Every $\mathbf{t}$ is a sum of rows that carry meaning on their own. The
prototypes are not constrained at all. They remain free vectors, but
they anchor a space that is now populated by attribute words, and a
prototype can therefore be read by comparing it to those words directly
in parameter space, with no detour through training data. Grounding the
prototypes is in this sense shorthand for grounding the space they
anchor, and the read-outs this buys are developed in
Section~\ref{sec:fi-readouts}.

\paragraph{The optional ID row.} The row $\mathbf{e}_{\mathrm{ID}(t)}$
in Equation~\ref{eq:fi-composition} is a free per-item vector added to
the sum by configuration. It restores per-item capacity, letting an item
deviate from what its attribute values alone would predict, at the price
of an ungrounded summand in every embedding. It also makes the host a
special case. With $\mathcal{F}_t$ empty and the ID row kept, the model
reduces to I-ProtoMF exactly, and our implementation reproduces the host
bit-identically in this configuration. The chapter therefore evaluates
three arms. fI-ProtoMF uses features and the ID row, the noid arm uses
features alone, and the f0 arm is the host reached through the extended
code path, the reference point of all comparisons. What the ID row
contributes, and what its ungrounded share costs, is measured in
Sections~\ref{sec:fi-readouts} and~\ref{sec:fi-results}.
